\section{Introduction}
\label{sec:introduction}

\PARstart{S}{olar} photovoltaic generation is now common in distribution-level energy planning, but its short-term behavior still creates practical trouble for control rooms, campus microgrids, and scheduling tools. The physical envelope of a PV plant is strongly constrained by sun position, season, and plant rating, while the usable inverter output can change quickly when clouds, humidity, wind, or rain alter the irradiance reaching the array. A forecasting model that ignores the deterministic solar component can waste model capacity relearning daily geometry. A model that relies only on a physical clear-sky curve misses the cloud-dependent departures that matter in operation.

This work studies 5-min PV power forecasting for a \PlantCapacity{} plant at \StudySite. The available data contain plant power, irradiance, theoretical power, temperature, humidity, wind speed, precipitation, total cloud cover, and low-, mid-, and high-cloud-cover layers. The study therefore sits between two common extremes. It is not a purely physical model, because the relationship between clouds and plant output is learned from data. It is also not a fully black-box model, because solar geometry, clear-sky behavior, and capacity constraints are explicit parts of the feature design.

The 5-min resolution is deliberately challenging. At this cadence, many errors are not caused by a failure to understand the daily solar arc; they are caused by local atmospheric transitions. A cloud edge may change output between adjacent samples, while an API-style weather forecast may report only cloud fraction, cloud layer, or shortwave radiation over a coarser spatial and temporal footprint. The forecasting task is not merely to learn a smooth bell-shaped daily curve. It must also handle days on which the same solar elevation produces very different power because the sky state has changed.

Baseline choice also matters in this setting. A previous-day same-time persistence forecast asks whether the next comparable clock time behaves like yesterday. A 5-min last-value persistence forecast asks a narrower question: given the plant output five minutes ago, what is the next value? The latter is a live nowcast with access to target-adjacent plant state. It is often hard to beat and should not be used to claim failure of a forecast model designed for weather-input deployment. This paper reports both baselines so the information setting is explicit.

The proposed method uses a two-stage cascade. Stage 1 estimates a clearness-index-like atmospheric transmissivity variable from weather, cloud, time, and solar-geometry inputs. Stage 2 predicts AC power using the Stage 1 output together with physical and operational features. Fig.~\ref{fig:architecture} summarizes the main idea. The design is intentionally practical: tabular inputs, XGBoost regressors, chronological validation, saved model artifacts, and a forecast-from-CSV inference path.

\begin{figure*}[!t]
\centering
\includegraphics[width=0.92\textwidth]{figures/diagrams/diagram_architecture.png}
\caption{Physics-guided two-stage forecasting architecture. The first model estimates an atmospheric transmissivity signal, while the second model combines that signal with solar-geometry, clear-sky, weather, and operational features to estimate 5-min AC power.}
\label{fig:architecture}
\end{figure*}

The contribution is methodological and application-oriented. The paper does not claim that XGBoost, clearness indices, or hybrid physical-data forecasting are new by themselves. Instead, the contribution is an integrated plant-level workflow that combines (i) high-resolution operational plant data, (ii) physics-guided feature engineering, (iii) a two-stage clearness-to-power cascade, (iv) chronological train-validation-test evaluation, (v) ablation against simpler model structures, and (vi) diagnostic figures that expose where the forecast is reliable and where weather-driven errors remain.

The remainder of the paper is organized as follows. Section~\ref{sec:related} positions the work against PV forecasting, energy-management, and physics-guided learning literature. Section~\ref{sec:dataset} describes the plant data and engineered variables. Section~\ref{sec:methodology} formulates the two-stage model and the main equations. Section~\ref{sec:experimental} gives the evaluation protocol. Section~\ref{sec:results} reports accuracy, ablation, and diagnostic behavior. Sections~\ref{sec:limitations}--\ref{sec:conclusion} discuss limits, future work, and conclusions.
